\subsection{Editorial evidence for Algebra 17136--18021}\label{algebra-proof-007-editorial-evidence-for-algebra-1713618021}

Authority: Stacks Project authors, algebra.tex, commit a04446e57ec1fbc252a871afcec7752fb2807b14, SHA-256 FA8BB92E58A4F78A2BD01B3B6A4A87DE0A0D279F5DD90641B574DD5FBFFFA4F3. This interval and its fifteen received reports have been read. Lines 18022--18030 are lookahead through the next statement and proof opening only. These arguments support the correction review; stronger results remain identifiable editorial material. No translated or cumulative chapter has been changed.

\subsubsection{Blowup fractions and base change}\label{algebra-proof-007-blowup-fractions-and-base-change}

Sources: definition-blow-up and lemma-affine-blowup, 17142--17181; lemma-blowup-base-change, 17183--17207; OCC-00424.

Retain the original graded algebra \(\bigoplus_{n\geq0}I^n\), including \(I^0=R\), and its degree-zero localization at \(a^{(1)}\). Its fractions \(x/a^n\), \(x\in I^n\), map to the same fractions in \(R_a\). The original equality relation \[
a^k(a^mx-a^ny)=0
\] is exactly the equality relation in \(R_a\), so this map is injective. This includes zero divisors in \(R\). The image is generated by all \(x/a\), \(x\in I\); products generate each \(I^n\). Multiplication by \(a\) on this subring of \(R_a\) is injective, \(IR'=aR'\) since \(x=a(x/a)\), and inverting \(a\) gives \(R_a\). For \(a=0\) the localization is the zero ring, where injectivity of its unique endomorphism is still valid.

For the base change retain \(b\), the image of \(a\), \(J=IS\), \(B=S\otimes_R R[I/a]\), and \(T_b(B)=\{u:\exists k\geq0,\ b^ku=0\}\). This is an ideal: a common larger exponent annihilates a sum, and multiplication by any element preserves annihilation. The map \[
\theta:B\longrightarrow S[J/b],\qquad
s\otimes(x/a^n)\longmapsto s\,\overline{x}/b^n
\] is well-defined by the original fraction equality, balancing, and the multiplication formula. It is surjective: every numerator in \(J^n=I^nS\) is a finite sum \(\sum_i\overline{x_i}s_i\), \(x_i\in I^n\). It kills \(T_b(B)\), because \(b\) is injective on the target.

There is also an exact direct proof of the claimed annihilation in \(S\otimes_R I^n\). If \(w=\sum_i s_i\otimes x_i\) maps to zero in \(J^n\), then \[
b^nw=\sum_i s_i\otimes a^nx_i
=\sum_i s_i\overline{x_i}\otimes a^n=0.
\] Here \(a^n\) lies in the original \(I^n\), so each balancing step has its specified domain. The source claim is valid in that tensor module, not merely after mapping into \(B\).

The inverse from \(S[J/b]\) to \(B/T_b(B)\) sends \(y/b^n\), with \(y=\sum_i\overline{x_i}s_i\), to \(\sum_i s_i\otimes x_i/a^n\) modulo torsion. A different decomposition at the same degree has difference killed by \(b^n\). For different representatives \(y/b^n=y'/b^m\), choose \(k\) with \(b^k(b^my-b^ny')=0\). If \(u,u'\) are the proposed images in \(B\), then \(b^nu=y\otimes1\) and \(b^mu'=y'\otimes1\), so \(b^{k+n+m}(u-u')=0\). The class is independent of both choices. Addition, multiplication and the identity follow by common denominators and these same balancing equations. On generators the two composites are the identity. Thus the original quotient statement holds for arbitrary ring maps, with every torsion term retained.

The report is a notation clarification: use \(b\)-power torsion and \(b^n\) on the \(S\)-side consistently. Identifying \(a\) with its image would also be legitimate if declared; it is not a false mathematical torsion assertion.

\subsubsection{Polynomial charts and the named base ring}\label{algebra-proof-007-polynomial-charts-and-the-named-base-ring}

Sources: examples at 17209--17247 and lemma-affine-blowup-quotient-description, 17249--17268; OCC-00422 and OCC-00423.

The source monomial is \(t_1^{e_1}\cdots t_n^{e_n}\), so the unrelated \(x_n\) is corrected to \(t_n\). The base-change map in the quotient-description proof is the declared \(\mathbf Z[t_1,\ldots,t_n]\to R\), not a new ring \(A\).

For the polynomial affine chart the displayed quotient identifies with \(R[t_1,x_2,\ldots,x_n]\) by the maps \(t_i\mapsto t_1x_i\) for \(i\geq2\), and the reverse map sending each retained variable to its class. The displayed relations verify both composites. Multiplication by \(t_1\) is injective even when \(R\) has zero divisors, by the monomial coefficients. After inverting \(t_1\), the maps \(x_i\leftrightarrow t_i/t_1\) are inverse. The source chart injects into its localization, as does the blowup fraction ring, so the surjection between them is injective. Base change along the stated map to \(R\), with the exact torsion quotient just proved, gives the claimed general presentation. Its \(n=1\) case retains the quotient of \(R\) by all \(a_1\)-power torsion; it need not equal \(R\).

\subsubsection{The field boundary and the directed chart construction}\label{algebra-proof-007-the-field-boundary-and-the-directed-chart-construction}

Source: lemma-valuation-ring-colimit-affine-blowups, 17351--17391; OCC-00425 and OCC-00426.

If \(R\) is a field, then \(\mathfrak m=0\) and \(R=A=K\). The original condition \(a\in I\subset\mathfrak m\) forces \(a=0\) and \(I=0\). That chart is zero and has zero fibre, so none of the prescribed objects exists. Preserve the theorem's field case by admitting the identity chart \((I,a)=(R,1)\), whose algebra is \(R\) and whose closed fibre is \(R/\mathfrak m\ne0\). Also state explicitly that the indexed charts are contained in the given \(A\); the properties of a chart alone do not select this particular valuation ring. These are corrections of the stated diagram, not a replacement of \(A\).

Suppose now \(R\) is not a field. For a finite set \(E=\{e_1,\ldots,e_n\}\subset A\), choose an actual common denominator \(0\ne d\in R\) with \(de_i\in R\). Choose \(0\ne t\in\mathfrak m\), and retain \(a=td\) and \(f_i=ae_i=t(de_i)\). Thus \(a,f_i\in\mathfrak m\) directly in \(R\), and \(a\ne0\). For \(I=(f_1,\ldots,f_n,a)\) the chart is \[
R[I/a]=R[e_1,\ldots,e_n]\subset A.
\] Indeed each fraction with numerator in \(I^k\) expands as a polynomial in the \(f_i/a=e_i\) and \(a/a=1\), and those generators are already chart fractions. If \(\mathfrak n\) is the maximal ideal of the dominating valuation ring, then \(\mathfrak n\cap R=\mathfrak m\); its contraction to the chart is a prime over \(\mathfrak m\). The chart's quotient by \(\mathfrak m\) is consequently nonzero, proving the required fibre condition.

For two such charts with pairs \((I,a),(J,b)\), their containing chart is \(R[IJ/(ab)]\). The first inclusion sends \(x/a^n\) to \(b^nx/(ab)^n\), retaining \(b^nx\in I^nJ^n\); the second inclusion is analogous. A numerator \(z\in(IJ)^n=I^nJ^n\) is a finite sum \(\sum_i x_i y_i\) and \[
z/(ab)^n=\sum_i(x_i/a^n)(y_i/b^n)\in A.
\] The product ideal remains finitely generated and contained in \(\mathfrak m\), the denominator remains nonzero, and contraction of \(\mathfrak n\) again gives a nonzero fibre. The identity chart combines with either chart to give that chart. Thus this is a nonempty directed system of subrings of the original \(K\), every finite subset of \(A\) lies in an object, and its colimit under the displayed inclusions is exactly \(A\). The empty finite subset is covered by the identity chart.

\subsubsection{The exact overring strengthening}\label{algebra-proof-007-the-exact-overring-strengthening}

The preceding argument does not require that the receiving overring be a valuation ring. For any local domain \((R,\mathfrak m)\) with fraction field \(K\), and any intermediate ring \(R\subset A\subset K\), the following are equivalent:

\begin{enumerate}
\def\labelenumi{\arabic{enumi}.}
\tightlist
\item
  \(A/\mathfrak mA\ne0\).
\item
  \(A\) is the directed union of its affine charts \(R[I/a]\subset A\) with \(I\) finitely generated, nonzero \(a\in I\subset\mathfrak m\), and nonzero closed fibre, together with the identity chart.
\end{enumerate}

For (1) choose a maximal ideal of \(A\) containing the proper ideal \(\mathfrak mA\). Its contraction to \(R\) is exactly \(\mathfrak m\). The finite-set common-denominator construction, product chart and fraction maps in the previous section apply without change; contraction of this ideal supplies the fibre condition for every chart inside \(A\). If \(R\) is a field, the intermediate ring is \(R\) and the identity chart suffices.

Conversely, if \(1=\sum_{j=1}^s r_j a_j\) with \(r_j\in\mathfrak m\), \(a_j\in A\), directedness puts the finite list \(a_j\) in one chart \(B\). Then \(1\in\mathfrak mB\), contrary to its nonzero closed fibre. Thus (2) implies (1). A finite-type \(A\) with nonzero closed fibre is itself one of the constructed charts, using a finite list of its algebra generators. This is a proved editorial strengthening; the translation retains the author's valuation-ring hypothesis and the visible field-case correction.

\subsubsection{Resolutions, chain maps and contravariance}\label{algebra-proof-007-resolutions-chain-maps-and-contravariance}

Sources: 17410--17754; OCC-00430, OCC-00427, OCC-01388, OCC-00428, OCC-00429, OCC-01389, OCC-01391 and OCC-01390.

In the finite-free construction the first kernel is \(\ker(R^{n_0}\to M)\). It is finite because \(R\) is Noetherian, so choose a finite free module surjecting onto it before using the printed recursion at \(e\geq1\). There is no \(R^{n_{-1}}\).

The chain identity is \[
\alpha_{i-1}d_{F,i}=d_{G,i}\alpha_i:F_i\longrightarrow G_{i-1}.
\] The original \(d_{G,i-1}\alpha_i\) is ill-typed. This correction is already present as MC-STK-ERR-1602.

For a chain map \(\alpha:F_\bullet\to G_\bullet\), the induced cochain map is \[
\alpha^*:\operatorname{Hom}_R(G_\bullet,N)
\longrightarrow\operatorname{Hom}_R(F_\bullet,N),\quad
u\longmapsto u\alpha_i\text{ in degree }i.
\] The original cochain differential is precomposition with \(d_{F,i+1}\); the chain identity verifies that this commutes with differentials. If \(\alpha-\beta=d_Gh+hd_F\), set \[
k^i(u)=u h_{i-1}.
\] Then \(\delta_F k^i(u)+k^{i+1}\delta_G(u)
=u h_{i-1}d_{F,i}+u d_{G,i+1}h_i=u(\alpha_i-\beta_i)\). Thus the cochain maps are homotopic and induce the same \(H^i\). In particular \(\alpha:F\to G\), \(\beta:G\to F\) give \[
H^i(\alpha^*)H^i(\beta^*)=H^i((\beta\alpha)^*)
\] on the cohomology of \(\operatorname{Hom}_R(F_\bullet,N)\), with identity \(\operatorname{id}_F\); the other composite is on \(G\). This verifies all parts of the retained MC-STK-ERR-1594 and MC-STK-ERR-1596 corrections. The \(H_i\) typo is already corrected by MC-STK-ERR-1595. The two grammar fixes are already present as MC-STK-ERR-1603 and MC-STK-ERR-1604. These six canonical units comprise nine operations and must not be applied a second time.

The remaining grammar changes at 17646 and 17698 leave both long exact sequences intact. Termwise free modules make Hom exact in the second variable; the first-variable construction is termwise split, so applying Hom reverses the arrows and remains exact. These are precisely the variance directions in the later annihilation and depth arguments.

\subsubsection{Projective comparison with the same maps}\label{algebra-proof-007-projective-comparison-with-the-same-maps}

The comparison lemma at 17499--17546 holds when each \(F_i\) is projective, rather than free, with no change to the arbitrary target resolution or the source complex condition. Here projective means that a map from \(F_i\) lifts through every surjection.

Let \(f:M\to N\) be the fixed map. Lift \(F_0\to M\to N\) through the surjection \(N_0\to N\) to obtain \(\alpha_0\). If \(\alpha_0,\ldots,\alpha_{i-1}\) are chosen, then \(\alpha_{i-1}d_{F,i}\) lands in \(\ker d_{N,i-1}\); for \(i=1\) this means the kernel of the augmentation. The assertion follows from the chain identities and \(d_F^2=0\), or from the augmentation condition when \(i=1\). Exactness makes \(N_i\to\ker d_{N,i-1}\) surjective, so projectivity gives \(\alpha_i\).

For two lifts put \(\gamma=\alpha-\beta\). Lift \(\gamma_0\) to \(h_0:F_0\to N_1\), as its augmentation is zero. Having chosen the lower \(h_j\), put \(u_i=\gamma_i-h_{i-1}d_{F,i}\). Then \[
d_{N,i}u_i
=\gamma_{i-1}d_{F,i}
-(\gamma_{i-1}-h_{i-2}d_{F,i-1})d_{F,i}=0.
\] Exactness and projectivity lift \(u_i\) to \(h_i:F_i\to N_{i+1}\). Hence \(\gamma_i=d_{N,i+1}h_i+h_{i-1}d_{F,i}\) at every degree. The augmentation supplies the initial case; no undefined negative module is needed.

Applying these constructions to projective resolutions in both directions yields inverse chain maps up to homotopy. The exact precomposition calculation above therefore identifies their Hom cohomology with the original Ext definition, naturally and contravariantly in the first module. Projective modules also make Hom exact in the second variable by the same lifting property. The original source chooses free resolutions, which already exist, and stays unchanged; this is supplementary generality, not a proposed replacement definition.

\subsubsection{Depth quotients and propagation to large powers}\label{algebra-proof-007-depth-quotients-and-propagation-to-large-powers}

Sources: 17759--18021; OCC-00431 and OCC-00432.

The source \(N\cong R/\mathfrak p\) gives the map \[
R/(\mathfrak p+(x))\longrightarrow N/xN,\qquad
r+(\mathfrak p+(x))\longmapsto r n_0+xN
\] for the chosen generator \(n_0\) with annihilator \(\mathfrak p\). It is surjective, and \(rn_0\in xN\) means \(r-xs\in\mathfrak p\) for some \(s\), proving its kernel is zero and giving the inverse through any representative. Thus both displayed occurrences require parentheses around \(\mathfrak p+(x)\). The original support, associated prime and dimension calculations keep the same quotient.

There is a uniform strengthening of lemma-inherit-minimal-primes. Let \(R\) be any Noetherian ring, \(M\) a finite module, \(\mathfrak p\in\operatorname{Ass}(M)\), and \(x\in R\). Choose the same \(N\cong R/\mathfrak p\subset M\). Artin--Rees gives \(c\geq0\) such that for every \(n\geq c\), \[
N\cap x^nM=x^{n-c}(N\cap x^cM).
\] For every \(n\geq c+1\), this is contained in \(xN\). Therefore \[
\overline N_n=N/(N\cap x^nM)\subset M/x^nM
\] surjects onto \(N/xN\), while it is annihilated by both \(\mathfrak p\) and \(x^n\). Consequently \[
\operatorname{Supp}(\overline N_n)=V(\mathfrak p+(x))
\] for all such \(n\): the surjection gives inclusion of the right-hand support, and the annihilators give the reverse inclusion. Every minimal prime of this closed set is an associated prime of the finite module \(\overline N_n\), hence of \(M/x^nM\). The single bound \(c+1\) works simultaneously for all primes minimal over \(\mathfrak p+(x)\). If that ideal is the unit ideal there are no such primes, and the assertion is empty. If \(x=0\), the formula remains valid and the assertion gives the original \(\mathfrak p\).

This proves the source conclusion for all sufficiently large \(n\), over an arbitrary Noetherian ring. It retains the actual submodule and its entire power-intersection formula. In the receiving depth proof at 17984--18020, choose any such \(n\); because its original \(x\) is a nonzerodivisor, so is \(x^n\), and the already established depth drop is one. No hypothesis of that local depth proof is removed.

The use of an \(x\) with depth drop exactly one in the earlier depth/Ext proof is not circular: finite depth attains a maximum length \(d\), choose the first element of a length-\(d\) regular sequence, and its tail shows depth at least \(d-1\). Any longer sequence on the quotient would concatenate with \(x\) to exceed \(d\), giving the reverse inequality. The Ext maps labelled by multiplication by this \(x\) vanish because \(x\) annihilates the original residue field. Induction gives finite first nonvanishing Ext and the stated equality. The later theorem then establishes this drop for every permitted nonzerodivisor. These source steps require no new correction.

\subsubsection{Reading boundaries and retained material}\label{algebra-proof-007-reading-boundaries-and-retained-material}

The entire current interval differs from the authority by nine already recorded canonical operations and eight reference/history environments attached to four Ext lemmas. Comparison after accounting for precisely those additions is exact. Those environments and all other existing work are preserved; comparison alone does not re-adjudicate their external FAC claims.

The existing corpus index was queried once for affine blowup. Its two hits were routed to original TeX. Bhatt--Scholze, Projectivity of the Witt vector affine Grassmannian, lines 251--284, supply contextual valuation-cover definitions and examples, with their cited dependencies not adopted here. Bhatt--Lurie, Absolute Prismatic Cohomology, lines 11899--11920, cover a regular-immersion blowup setup and only the start of the syntomic formula, not its proof. Bezrukavnikov--Finkelberg--Mirkovic, Equivariant (K-)homology of affine Grassmannian and Toda lattice, lines 1695--1814, cover the earlier symbolic-chart proposition through its proof, comments and the next proof opening. Its root-theoretic and flatness dependencies are not fully reviewed, and it is not adopted as a proof of the present algebraic claims. In particular symbolic powers in that setting are not substituted for the ordinary Rees powers here.

The proofs used for this batch are the exact supplied Stacks arguments and the explicit comparisons above. No exhaustive reading, novelty claim, global synthesis, rendering or publication is claimed. Broader synthesis follows the core report review.
